%=================================================================
% LuaLaTeX preamble
% (%=================================================================
% LuaLaTeX preamble
% (%=================================================================
% LuaLaTeX preamble
% (%=================================================================
% LuaLaTeX preamble
% (\input{preamble}でtemplate.texに読み込まれることを想定)
%
% 用途:
%   日本語・英語・古典ギリシア語混在文書
%   (哲学・神学・古典学向け)
%
% 2026-06
%
% 重要:
%   - Unicode Greek を使用可能
%   - Babel/LGR の古い \textgreek{...} 入力を維持
%   - hyperref による PDF しおり対応
%=================================================================

%=================================================================
% 日本語処理
%=================================================================

\usepackage{luatexja-fontspec}
\usepackage{luatexja-ruby}

% 古典ギリシア語(Unicode)を和文扱いしない。
%
% これが無いと
%
%   λόγος
%   ἀλήθεια
%
% などが HaranoAjiMincho に送られ、
%
%   Missing character: There is no ό ...
%
% が発生する。
%
\ltjsetparameter{jacharrange={-2,-3}}

%=================================================================
% フォント
%=================================================================

% 欧文・ギリシア語
\setmainfont{Libertinus Serif}

% 和文
\setmainjfont{HaranoAjiMincho}

%=================================================================
% 古典ギリシア語
%=================================================================

% Babel の古い LGR 入力
%
%   \textgreek{l'ogos}
%   \textgreek{>'anjrwpoc}
%
% を維持するために必要。
%
% Unicode Greek
%
%   λόγος
%   \foreignlanguage{greek}{πάντες ἄνθρωποι}
%
% と共存可能。
%
%  しかしitalic, boldなどが出なくなるので、fontencは外す。その結果、LGRは不可
% \usepackage[LGR,T1]{fontenc}
\usepackage[polutonikogreek,english,japanese]{babel}

%=================================================================
% レイアウト
%=================================================================

\usepackage{paracol}
\footnotelayout{m}

\usepackage{longtable}

%=================================================================
% 色
%=================================================================

\usepackage{xcolor}

%=================================================================
% PDF ハイパーリンク
%=================================================================

\usepackage[
unicode,
bookmarks=true,
hidelinks
]{hyperref}

% hyperref の後に読むこと。
% 先に読むと Option clash が発生する。
%
\usepackage{bookmark}

%=================================================================
% ヘッダ・フッタ
%=================================================================

\usepackage{fancyhdr}
\pagestyle{fancy}
%=================================================================
% 注意事項
%=================================================================
%
% 見出し中のギリシア語:
%
%   \section{λόγος}
%
% は空白になる場合がある。
%
% その場合:
%
%   \section{{\rmfamily λόγος}}
%
% と書く。
%
%
% Babel ギリシア語:
%
%   \textgreek{l'ogos}
%
% Unicode ギリシア語:
%
%   \foreignlanguage{greek}{λόγος}
%
% の両方が利用可能。
%
%=================================================================
でtemplate.texに読み込まれることを想定)
%
% 用途:
%   日本語・英語・古典ギリシア語混在文書
%   (哲学・神学・古典学向け)
%
% 2026-06
%
% 重要:
%   - Unicode Greek を使用可能
%   - Babel/LGR の古い \textgreek{...} 入力を維持
%   - hyperref による PDF しおり対応
%=================================================================

%=================================================================
% 日本語処理
%=================================================================

\usepackage{luatexja-fontspec}
\usepackage{luatexja-ruby}

% 古典ギリシア語(Unicode)を和文扱いしない。
%
% これが無いと
%
%   λόγος
%   ἀλήθεια
%
% などが HaranoAjiMincho に送られ、
%
%   Missing character: There is no ό ...
%
% が発生する。
%
\ltjsetparameter{jacharrange={-2,-3}}

%=================================================================
% フォント
%=================================================================

% 欧文・ギリシア語
\setmainfont{Libertinus Serif}

% 和文
\setmainjfont{HaranoAjiMincho}

%=================================================================
% 古典ギリシア語
%=================================================================

% Babel の古い LGR 入力
%
%   \textgreek{l'ogos}
%   \textgreek{>'anjrwpoc}
%
% を維持するために必要。
%
% Unicode Greek
%
%   λόγος
%   \foreignlanguage{greek}{πάντες ἄνθρωποι}
%
% と共存可能。
%
%  しかしitalic, boldなどが出なくなるので、fontencは外す。その結果、LGRは不可
% \usepackage[LGR,T1]{fontenc}
\usepackage[polutonikogreek,english,japanese]{babel}

%=================================================================
% レイアウト
%=================================================================

\usepackage{paracol}
\footnotelayout{m}

\usepackage{longtable}

%=================================================================
% 色
%=================================================================

\usepackage{xcolor}

%=================================================================
% PDF ハイパーリンク
%=================================================================

\usepackage[
unicode,
bookmarks=true,
hidelinks
]{hyperref}

% hyperref の後に読むこと。
% 先に読むと Option clash が発生する。
%
\usepackage{bookmark}

%=================================================================
% ヘッダ・フッタ
%=================================================================

\usepackage{fancyhdr}
\pagestyle{fancy}
%=================================================================
% 注意事項
%=================================================================
%
% 見出し中のギリシア語:
%
%   \section{λόγος}
%
% は空白になる場合がある。
%
% その場合:
%
%   \section{{\rmfamily λόγος}}
%
% と書く。
%
%
% Babel ギリシア語:
%
%   \textgreek{l'ogos}
%
% Unicode ギリシア語:
%
%   \foreignlanguage{greek}{λόγος}
%
% の両方が利用可能。
%
%=================================================================
でtemplate.texに読み込まれることを想定)
%
% 用途:
%   日本語・英語・古典ギリシア語混在文書
%   (哲学・神学・古典学向け)
%
% 2026-06
%
% 重要:
%   - Unicode Greek を使用可能
%   - Babel/LGR の古い \textgreek{...} 入力を維持
%   - hyperref による PDF しおり対応
%=================================================================

%=================================================================
% 日本語処理
%=================================================================

\usepackage{luatexja-fontspec}
\usepackage{luatexja-ruby}

% 古典ギリシア語(Unicode)を和文扱いしない。
%
% これが無いと
%
%   λόγος
%   ἀλήθεια
%
% などが HaranoAjiMincho に送られ、
%
%   Missing character: There is no ό ...
%
% が発生する。
%
\ltjsetparameter{jacharrange={-2,-3}}

%=================================================================
% フォント
%=================================================================

% 欧文・ギリシア語
\setmainfont{Libertinus Serif}

% 和文
\setmainjfont{HaranoAjiMincho}

%=================================================================
% 古典ギリシア語
%=================================================================

% Babel の古い LGR 入力
%
%   \textgreek{l'ogos}
%   \textgreek{>'anjrwpoc}
%
% を維持するために必要。
%
% Unicode Greek
%
%   λόγος
%   \foreignlanguage{greek}{πάντες ἄνθρωποι}
%
% と共存可能。
%
%  しかしitalic, boldなどが出なくなるので、fontencは外す。その結果、LGRは不可
% \usepackage[LGR,T1]{fontenc}
\usepackage[polutonikogreek,english,japanese]{babel}

%=================================================================
% レイアウト
%=================================================================

\usepackage{paracol}
\footnotelayout{m}

\usepackage{longtable}

%=================================================================
% 色
%=================================================================

\usepackage{xcolor}

%=================================================================
% PDF ハイパーリンク
%=================================================================

\usepackage[
unicode,
bookmarks=true,
hidelinks
]{hyperref}

% hyperref の後に読むこと。
% 先に読むと Option clash が発生する。
%
\usepackage{bookmark}

%=================================================================
% ヘッダ・フッタ
%=================================================================

\usepackage{fancyhdr}
\pagestyle{fancy}
%=================================================================
% 注意事項
%=================================================================
%
% 見出し中のギリシア語:
%
%   \section{λόγος}
%
% は空白になる場合がある。
%
% その場合:
%
%   \section{{\rmfamily λόγος}}
%
% と書く。
%
%
% Babel ギリシア語:
%
%   \textgreek{l'ogos}
%
% Unicode ギリシア語:
%
%   \foreignlanguage{greek}{λόγος}
%
% の両方が利用可能。
%
%=================================================================
でtemplate.texに読み込まれることを想定)
%
% 用途:
%   日本語・英語・古典ギリシア語混在文書
%   (哲学・神学・古典学向け)
%
% 2026-06
%
% 重要:
%   - Unicode Greek を使用可能
%   - Babel/LGR の古い \textgreek{...} 入力を維持
%   - hyperref による PDF しおり対応
%=================================================================

%=================================================================
% 日本語処理
%=================================================================

\usepackage{luatexja-fontspec}
\usepackage{luatexja-ruby}

% 古典ギリシア語(Unicode)を和文扱いしない。
%
% これが無いと
%
%   λόγος
%   ἀλήθεια
%
% などが HaranoAjiMincho に送られ、
%
%   Missing character: There is no ό ...
%
% が発生する。
%
\ltjsetparameter{jacharrange={-2,-3}}

%=================================================================
% フォント
%=================================================================

% 欧文・ギリシア語
\setmainfont{Libertinus Serif}

% 和文
\setmainjfont{HaranoAjiMincho}

%=================================================================
% 古典ギリシア語
%=================================================================

% Babel の古い LGR 入力
%
%   \textgreek{l'ogos}
%   \textgreek{>'anjrwpoc}
%
% を維持するために必要。
%
% Unicode Greek
%
%   λόγος
%   \foreignlanguage{greek}{πάντες ἄνθρωποι}
%
% と共存可能。
%
%  しかしitalic, boldなどが出なくなるので、fontencは外す。その結果、LGRは不可
% \usepackage[LGR,T1]{fontenc}
\usepackage[polutonikogreek,english,japanese]{babel}

%=================================================================
% レイアウト
%=================================================================

\usepackage{paracol}
\footnotelayout{m}

\usepackage{longtable}

%=================================================================
% 色
%=================================================================

\usepackage{xcolor}

%=================================================================
% PDF ハイパーリンク
%=================================================================

\usepackage[
unicode,
bookmarks=true,
hidelinks
]{hyperref}

% hyperref の後に読むこと。
% 先に読むと Option clash が発生する。
%
\usepackage{bookmark}

%=================================================================
% ヘッダ・フッタ
%=================================================================

\usepackage{fancyhdr}
\pagestyle{fancy}
%=================================================================
% 注意事項
%=================================================================
%
% 見出し中のギリシア語:
%
%   \section{λόγος}
%
% は空白になる場合がある。
%
% その場合:
%
%   \section{{\rmfamily λόγος}}
%
% と書く。
%
%
% Babel ギリシア語:
%
%   \textgreek{l'ogos}
%
% Unicode ギリシア語:
%
%   \foreignlanguage{greek}{λόγος}
%
% の両方が利用可能。
%
%=================================================================
